% Part: cheat-sheets
% Chapter: many-valued-logic

\documentclass[../../../../include/open-logic-chapter]{subfiles}

\begin{document}

\olchapter{chs}{mvl}{Many-valued Logic}

\section{Semantics}

\begin{defn}[Matrix]
\ollabel{defn:matrix}
A \emph{matrix} for the logic~$\Log L$ consists of:
\begin{enumerate}
\item a set of connectives making up a language~$\Lang L$;
\item a set $V \neq \emptyset$ of truth values;
\item a set $V^+ \subseteq V$ of designated truth values;
\item for each $n$-place connective $\star$ in $\Lang L$, a truth
function~$\tf{\star} : V^n \to V$. If $n = 0$, then $\tf{\star}$ is
just an element of~$V$.
\end{enumerate}
\end{defn}

\begin{ex}
The matrix for classical logic~\LogCL{} consists of:
\begin{enumerate}
  \item The standard propositional language $\Lang L_0$ with
  $\lfalse$, $\lnot$, $\land$, $\lor$, $\lif$.
  \item The set of truth values $V = \{\True, \False\}$.
  \item $\True$ is the only designated value, i.e., $V^+ = \{\True\}$.
  \item For $\lfalse$, we have $\tf{\lfalse} = \False$. The other
  truth functions are given by the usual truth tables (see
  \olref{fig:tf-CL}).
\end{enumerate}
\begin{figure}
  \begin{center}
      \begin{tabular}{c|c} 
        $\tf{\lnot}$ & \\ 
        \hline  
        $\True$ & $\False$ \\ 
        $\False$ & $\True$ 
      \end{tabular}
      \quad
      \begin{tabular}{c|cc} 
        $\tf{\land}$ & $\True$ & $\False$ \\ 
        \hline 
        $\True$ & $\True$ & $\False$ \\ 
        $\False$ & $\False$ & $\False$ 
      \end{tabular}
      \quad
      \begin{tabular}{c|cc} 
        $\tf{\lor}$ & $\True$ & $\False$ \\ 
        \hline 
        $\True$ & $\True$ & $\True$ \\ 
        $\False$ & $\True$ & $\False$ 
      \end{tabular}
      \quad
      \begin{tabular}{c|cc} 
        $\tf{\lif}$ & $\True$ & $\False$ \\ 
        \hline 
        $\True$ & $\True$ & $\False$ \\ 
        $\False$ & $\True$ & $\True$ 
      \end{tabular}
    \end{center} 
    \caption{Truth functions for classical logic~$\LogCL$.}
    \ollabel{fig:tf-CL}
  \end{figure}
  \end{ex}

\begin{defn}\ollabel{def:lukasiewicz}
Three-valued \L ukasiewicz logic is defined using the matrix:
\begin{enumerate}
  \item The standard propositional language $\Lang L_0$ with
  $\lnot$, $\land$, $\lor$, $\lif$.
  \item The set of truth values $V = \{\True, \Undef, \False\}$.
  \item $\True$ is the only designated value, i.e., $V^+ = \{\True\}$.
  \item Truth functions are given by the following tables:
  \begin{center}
    \begin{tabular}{c|c} 
      $\tf{\lnot}$ & \\ 
      \hline  
      $\True$ & $\False$ \\ 
      $\Undef$ & $\Undef$ \\
      $\False$ & $\True$ 
    \end{tabular}
    \quad
    \begin{tabular}{c|ccc} 
      $\tf{\land}[\LogLuk[3]]$ & $\True$ & $\Undef$ & $\False$ \\ 
      \hline 
      $\True$ & $\True$ & $\Undef$ & $\False$ \\ 
      $\Undef$ & $\Undef$ & $\Undef$ & $\False$\\ 
      $\False$ & $\False$ & $\False$ & $\False$ 
    \end{tabular}
    \\[2ex]
    \begin{tabular}{c|ccc} 
      $\tf{\lor}[\LogLuk[3]]$ & $\True$ & $\Undef$ & $\False$ \\ 
      \hline 
      $\True$ & $\True$ & $\True$ & $\True$ \\ 
      $\Undef$ & $\True$ & $\Undef$ & $\Undef$ \\
      $\False$ & $\True$ & $\Undef$ & $\False$ 
    \end{tabular}
    \quad
    \begin{tabular}{c|ccc} 
      $\tf{\lif}[\LogLuk[3]]$ & $\True$ & $\Undef$ & $\False$ \\ 
      \hline 
      $\True$ & $\True$ & $\Undef$ & $\False$ \\ 
      $\Undef$ & $\True$ & $\True$ & $\Undef$  \\ 
      $\False$ & $\True$ & $\True$ & $\True$ 
    \end{tabular}
  \end{center} 
\end{enumerate}
\end{defn}

\begin{defn}
\emph{Strong Kleene logic}~$\LogKs$ is defined using the matrix:
\begin{enumerate}
  \item The standard propositional language $\Lang L_0$ with
  $\lnot$, $\land$, $\lor$, $\lif$.
  \item The set of truth values $V = \{\True, \Undef, \False\}$.
  \item $\True$ is the only designated value, i.e., $V^+ = \{\True\}$.
  \item Truth functions are given by the following tables:
  \begin{center}
    \begin{tabular}{c|c} 
      $\tf{\lnot}$ & \\ 
      \hline  
      $\True$ & $\False$ \\ 
      $\Undef$ & $\Undef$ \\
      $\False$ & $\True$ 
    \end{tabular}
    \quad
    \begin{tabular}{c|ccc} 
      $\tf{\land}[\LogKs]$ & $\True$ & $\Undef$ & $\False$ \\ 
      \hline 
      $\True$ & $\True$ & $\Undef$ & $\False$ \\ 
      $\Undef$ & $\Undef$ & $\Undef$ & $\False$\\ 
      $\False$ & $\False$ & $\False$ & $\False$ 
    \end{tabular}
    \\[2ex]
    \begin{tabular}{c|ccc} 
      $\tf{\lor}[\LogKs]$ & $\True$ & $\Undef$ & $\False$ \\ 
      \hline 
      $\True$ & $\True$ & $\True$ & $\True$ \\ 
      $\Undef$ & $\True$ & $\Undef$ & $\Undef$ \\
      $\False$ & $\True$ & $\Undef$ & $\False$ 
    \end{tabular}
    \quad
    \begin{tabular}{c|ccc} 
      $\tf{\lif}[\LogKs]$ & $\True$ & $\Undef$ & $\False$ \\ 
      \hline 
      $\True$ & $\True$ & $\Undef$ & $\False$ \\ 
      $\Undef$ & $\True$ & $\Undef$ & $\Undef$  \\ 
      $\False$ & $\True$ & $\True$ & $\True$ 
    \end{tabular}
  \end{center} 
\end{enumerate}
\end{defn}

\begin{defn}
\emph{Weak Kleene logic}~$\LogKw$ is defined using the matrix:
\begin{enumerate}
  \item The standard propositional language $\Lang L_0$ with
  $\lnot$, $\land$, $\lor$, $\lif$.
  \item The set of truth values $V = \{\True, \Undef, \False\}$.
  \item $\True$ is the only designated value, i.e., $V^+ = \{\True\}$.
  \item Truth functions are given by the following tables:
  \begin{center}
    \begin{tabular}{c|c} 
      $\tf{\lnot}$ & \\ 
      \hline  
      $\True$ & $\False$ \\ 
      $\Undef$ & $\Undef$ \\
      $\False$ & $\True$ 
    \end{tabular}
    \quad
    \begin{tabular}{c|ccc} 
      $\tf{\land}[\LogKw]$ & $\True$ & $\Undef$ & $\False$ \\ 
      \hline 
      $\True$ & $\True$ & $\Undef$ & $\False$ \\ 
      $\Undef$ & $\Undef$ & $\Undef$ & $\Undef$\\ 
      $\False$ & $\False$ & $\Undef$ & $\False$ 
    \end{tabular}
    \\[2ex]
    \begin{tabular}{c|ccc} 
      $\tf{\lor}[\LogKw]$ & $\True$ & $\Undef$ & $\False$ \\ 
      \hline 
      $\True$ & $\True$ & $\Undef$ & $\True$ \\ 
      $\Undef$ & $\Undef$ & $\Undef$ & $\Undef$ \\
      $\False$ & $\True$ & $\Undef$ & $\False$ 
    \end{tabular}
    \quad
    \begin{tabular}{c|ccc} 
      $\tf{\lif}[\LogKw]$ & $\True$ & $\Undef$ & $\False$ \\ 
      \hline 
      $\True$ & $\True$ & $\Undef$ & $\False$ \\ 
      $\Undef$ & $\Undef$ & $\Undef$ & $\Undef$  \\ 
      $\False$ & $\True$ & $\Undef$ & $\True$ 
    \end{tabular}
  \end{center} 
\end{enumerate}
\end{defn}

\begin{defn}
The \emph{logic of paradox}~$\LogLP$ is defined using the matrix:
\begin{enumerate}
  \item The standard propositional language $\Lang L_0$ with
  $\lnot$, $\land$, $\lor$, $\lif$.
  \item The set of truth values $V = \{\True, \Undef, \False\}$.
  \item $\True$ and $\Undef$ are designated, i.e., $V^+ = \{\True, \Undef\}$.
  \item Truth functions are the same as in strong Kleene logic.
\end{enumerate}
\end{defn}

\section{Tableaux}

\subsection{Hypotheses to test for validity}

To test $\Gamma \Entails[\Log[X]] !A$, suppose that $\pValue{!B_i} \in V^+$ for any 
$!B_i\in\Gamma$ and $\pValue{!A} \notin V^+$. 

\begin{enumerate}
	\item If $V^+ = \{\True\}$ ($\LogLuk[3]$, $\LogKs$, $\LogKw$), assume that
	the premises are $\True$ and the conclusion is either $\Undef$ or $\False$ (two branches).
	\item If $V^+ = \{\True, \Undef\}$, assume that the conclusion is $\False$ and for each 
	premise $!B_i$, divide every branch into one for $\sFmla{\True}{!B_i}$ and one for
	 $sFmla{\False}{!B_i}$. (E.g. with 2 premises, 4 branches, and with 3, 27.)
\end{enumerate}

\subsection{Rules for $\lnot$}

In any of the logics $\LogLuk[3]$, $\LogKs$, $\LogKw$, $\LogLP$:

\begin{defish}
\AxiomC{\sFmla{\True}{\lnot !A}}
\RightLabel{\TRule{\True}{\lnot}}
\UnaryInfC{\sFmla{\False}{!A}}
\DisplayProof
\hfill
\AxiomC{\sFmla{\Undef}{\lnot !A}}
\RightLabel{\TRule{\Undef}{\lnot}}
\UnaryInfC{\sFmla{\Undef}{!A}}
\DisplayProof
\hfill
\AxiomC{\sFmla{\False}{\lnot !A}}
\RightLabel{\TRule{\False}{\lnot}}
\UnaryInfC{\sFmla{\True}{!A}}
\DisplayProof
\end{defish}

\subsection{Rules for $\land$ and $\lor$}

\subsubsection{Logics $\LogLuk[3]$, $\LogKs$, $\LogLP$}

Rules for ($\land$):

\begin{defish}\noindent
\AxiomC{\sFmla{\True}{!A \land !B}}
\RightLabel{\TRule{\True}{\land}}
\UnaryInfC{\sFmla{\True}{!A}}
\noLine
\UnaryInfC{\sFmla{\True}{!B}}
\DisplayProof
\hfill
\AxiomC{\sFmla{\Undef}{!A \land !B}}
\RightLabel{\TRule{\Undef}{\land}}
\UnaryInfC{$\sFmla{\True}{!A} \quad \mid \quad \sFmla{\Undef}{!A} \quad \mid \quad \sFmla{\Undef}{!A}$}
\noLine
\UnaryInfC{$\sFmla{\Undef}{!B} \quad \mid \quad \sFmla{\True}{!B} \quad \mid \quad \sFmla{\Undef}{!B}$}
\DisplayProof
\hfill
\AxiomC{\sFmla{\False}{!A \land !B}}
\RightLabel{\TRule{\False}{\land}}
\UnaryInfC{$\sFmla{\False}{!A} \quad \mid \quad \sFmla{\False}{!B}$}
\DisplayProof
\end{defish}

Rules for $\lor$:

\begin{defish}\noindent
\AxiomC{\sFmla{\True}{!A \lor !B}}
\RightLabel{\TRule{\True}{\lor}}
\UnaryInfC{$\sFmla{\True}{!A} \quad \mid \quad \sFmla{\True}{!B}$}
\DisplayProof
\hfill
\AxiomC{\sFmla{\Undef}{!A \lor !B}}
\RightLabel{\TRule{\Undef}{\lor}}
\UnaryInfC{$\sFmla{\False}{!A} \quad \mid \quad \sFmla{\Undef}{!A} \quad \mid \quad \sFmla{\Undef}{!A}$}
\noLine
\UnaryInfC{$\sFmla{\Undef}{!B} \quad \mid \quad \sFmla{\False}{!B} \quad \mid \quad \sFmla{\Undef}{!B}$}
\DisplayProof
\hfill
\AxiomC{\sFmla{\False}{!A \lor !B}}
\RightLabel{\TRule{\False}{\lor}}
\UnaryInfC{\sFmla{\False}{!A}}
\noLine
\UnaryInfC{\sFmla{\False}{!B}}
\DisplayProof
\end{defish}

\subsubsection{Logic $\LogKw$}

Rules for ($\land$):

\begin{defish}\noindent
\AxiomC{\sFmla{\True}{!A \land !B}}
\RightLabel{\TRule{\True}{\land}}
\UnaryInfC{\sFmla{\True}{!A}}
\noLine
\UnaryInfC{\sFmla{\True}{!B}}
\DisplayProof
\hfill
\AxiomC{\sFmla{\Undef}{!A \land !B}}
\RightLabel{\TRule{\Undef}{\land}}
\UnaryInfC{$\sFmla{\Undef}{!A} \quad \mid \quad \sFmla{\Undef}{!B}$}
\DisplayProof
\hfill
\AxiomC{\sFmla{\False}{!A \land !B}}
\RightLabel{\TRule{\False}{\land}}
\UnaryInfC{$\sFmla{\True}{!A} \quad \mid \quad \sFmla{\False}{!B} \quad \mid \quad \sFmla{\False}{!A}$}
\noLine
\UnaryInfC{$\sFmla{\False}{!B} \quad \mid \quad \sFmla{\True}{!A}\quad \mid \quad \sFmla{\False}{!B} $}
\DisplayProof
\end{defish}

Rules for $\lor$:

\begin{defish}\noindent
\AxiomC{\sFmla{\True}{!A \lor !B}}
\RightLabel{\TRule{\True}{\lor}}
\UnaryInfC{$\sFmla{\True}{!A} \quad \mid \quad \sFmla{\True}{!B} \quad \mid \quad \sFmla{\False}{!A}$}
\noLine
\UnaryInfC{$\sFmla{\True}{!B} \quad \mid \quad \sFmla{\False}{!A}\quad \mid \quad \sFmla{\True}{!B} $}
\DisplayProof
\hfill
\AxiomC{\sFmla{\Undef}{!A \lor !B}}
\RightLabel{\TRule{\Undef}{\lor}}
\UnaryInfC{$\sFmla{\Undef}{!A} \quad \mid \quad \sFmla{\Undef}{!B}$}
\DisplayProof
\hfill
\AxiomC{\sFmla{\False}{!A \lor !B}}
\RightLabel{\TRule{\False}{\lor}}
\UnaryInfC{\sFmla{\False}{!A}}
\noLine
\UnaryInfC{\sFmla{\False}{!B}}
\DisplayProof
\end{defish}

\subsection{Rules for $\lif$}

\subsubsection{Logics $\LogKw$}

Easy, left as an exercise.

\subsubsection{Logics $\LogLuk[3]$, $\LogKs$, $\LogLP$}

In $\LogKs$ (and $\LogLP$, which uses the same table) the conditional $!A \lif !B$ is equivalent to $\lnot !A \lor B$. 
The rules can be inferred from those for the disjunction.

\begin{defish}\noindent
\AxiomC{\sFmla{\True}{!A \lif !B}}
\RightLabel{\TRule{\True}{\lif}}
\UnaryInfC{$\sFmla{\False}{!A} \quad \mid \quad \sFmla{\True}{!B}$}
\DisplayProof
\hfill
\AxiomC{\sFmla{\Undef}{!A \lif !B}}
\RightLabel{\TRule{\Undef}{\lif}}
\UnaryInfC{$\sFmla{\True}{!A} \quad \mid \quad \sFmla{\Undef}{!A} \quad \mid \quad \sFmla{\Undef}{!A}$}
\noLine
\UnaryInfC{$\sFmla{\Undef}{!B} \quad \mid \quad \sFmla{\False}{!B} \quad \mid \quad \sFmla{\Undef}{!B}$}
\DisplayProof
\hfill
\AxiomC{\sFmla{\False}{!A \lif !B}}
\RightLabel{\TRule{\False}{\lif}}
\UnaryInfC{\sFmla{\True}{!A}}
\noLine
\UnaryInfC{\sFmla{\False}{!B}}
\DisplayProof
\end{defish}

\subsubsection{Rules for $\lif$ in $\LogLuk[3]$}

In $\LogLuk[3]$, there is one more way for the conditional to be true, and one less for it to be $\Undef$. 

\begin{defish}\noindent
\AxiomC{\sFmla{\True}{!A \lif !B}}
\RightLabel{\TRule{\True}{\lif}}
\UnaryInfC{$\sFmla{\Undef}{!A} \quad \mid \quad \sFmla{\False}{!A} \quad \mid \quad \sFmla{\True}{!B}$}
\noLine
\UnaryInfC{$\sFmla{\Undef}{!B} \quad \mid \quad \sFmla{~~~}{~~~} \quad \mid \quad \sFmla{~~~}{~~~}$}
\DisplayProof
\hfill
\AxiomC{\sFmla{\Undef}{!A \lif !B}}
\RightLabel{\TRule{\Undef}{\lif}}
\UnaryInfC{$\sFmla{\True}{!A} \quad \mid \quad \sFmla{\Undef}{!A}$}
\noLine
\UnaryInfC{$\sFmla{\Undef}{!B} \quad \mid \quad \sFmla{\False}{!B}$}
\DisplayProof
\hfill
\AxiomC{\sFmla{\False}{!A \lif !B}}
\RightLabel{\TRule{\False}{\lif}}
\UnaryInfC{\sFmla{\True}{!A}}
\noLine
\UnaryInfC{\sFmla{\False}{!B}}
\DisplayProof
\end{defish}

\OLEndChapterHook

\end{document}
